\subsection{Scope and Report Organization}

\paragraph{Objectives.}
The objective of this project is to build a \emph{decision support system} (DSS) for hotel revenue management, utilising a data warehouse and data analytics as required by the course specification. The system supports a revenue manager in making three recurring tactical decisions --- pricing, group-versus-transient allocation, and overbooking --- defined precisely in Section~2. For each decision, the system produces measurable outputs aligned with the three key performance indicators of the lodging industry~\cite{goeldner2009} --- ADR, occupancy, and RevPAR. The application is interactive: the user can adjust model parameters (price levels, group allocation, overbooking thresholds) and immediately observe the resulting changes in projected KPIs through \emph{what-if analysis}, \emph{sensitivity analysis}, and \emph{goal seeking} --- the three analysis modes that Turban et al.~identify as core DSS capabilities~\cite{turban2011}.

\paragraph{Scope.}
In scope: a dimensional data warehouse built from the Antonio et al.~dataset~\cite{antonio2019}; a model base of five declared models managed by a model base management system (MBMS) --- the core of the model management subsystem (the course-adopted topic); and an interactive web application exposing the models to the decision maker. Out of scope, for reasons established in this section: distribution-channel management and competitor pricing (the dataset carries no channel commissions or competitor rates --- pricing optimisation's ``competitor pricing'' input~\cite{amadeus2023} cannot be supported); function-space, restaurant, and ancillary revenue (room revenue only~\cite{goeldner2009}); and real-time operation (the dataset is a static historical research extract rather than daily-refreshed data~\cite{amadeus2023}).

\paragraph{Report organization.}
The remainder of the report is organised as follows. Section~2 defines the business problem: the organisational context, the decision makers, the three recurring decisions and their classification, the data, and the expected goals. Section~3 presents the proposed solution: the DSS architecture, the star schema of the data warehouse, the model catalog, and the model management subsystem presented in depth. Section~4 describes the application development, Section~5 evaluates the solution, and Section~6 concludes the report.
